% !TEX TS-program = xelatex
% !TEX encoding = UTF-8 Unicode
% !TEX TS-program = xelatex
% !TEX encoding = UTF-8 Unicode
\documentclass[aspectratio=169,11pt]{beamer}

\usetheme{Madrid}
\usecolortheme{default}
\setbeamertemplate{navigation symbols}{}
\setbeamertemplate{caption}[numbered]

\usepackage{fontspec}
\setsansfont{Latin Modern Sans}
\setmainfont{Latin Modern Roman}
\setmonofont{Latin Modern Mono}
\usepackage{graphicx}
\usepackage{booktabs}
\usepackage{array}
\usepackage{tabularx}
\usepackage{amsmath}

\definecolor{accentblue}{RGB}{22,58,105}
\definecolor{accentteal}{RGB}{20,122,112}
\definecolor{accentorange}{RGB}{232,116,28}
\definecolor{softblue}{RGB}{238,246,255}
\definecolor{ink}{RGB}{25,33,45}

\setbeamercolor{title}{fg=white}
\setbeamercolor{subtitle}{fg=white}
\setbeamercolor{frametitle}{fg=white,bg=accentblue}
\setbeamercolor{structure}{fg=accentblue}
\setbeamercolor{block title}{fg=white,bg=accentblue}
\setbeamercolor{block body}{fg=ink,bg=softblue}

\newcommand{\reportclaim}[1]{%
  \vspace{0.08cm}
  {\large\bfseries\color{accentblue}#1\par}
  \vspace{0.06cm}
}

\newcommand{\smallnote}[1]{%
  \par
  \vspace{0.08cm}
  {\footnotesize\color{ink}#1\par}
}

\newcommand{\widefig}[2][0.66\textheight]{%
  \makebox[\textwidth][c]{\includegraphics[width=0.96\textwidth,height=#1,keepaspectratio]{#2}}%
}

\newcommand{\chapterdivider}[2]{%
  \begin{frame}[plain]
    \centering
    \vspace{1.2cm}
    \makebox[\textwidth][c]{%
      \begin{beamercolorbox}[wd=0.84\textwidth,rounded=true,sep=0.42cm,center]{frametitle}
        {\LARGE\bfseries #1\par}
        \vspace{0.20cm}
        {\large #2\par}
      \end{beamercolorbox}%
    }
  \end{frame}
}

\newcommand{\twopic}[4]{%
  \begin{columns}[T,totalwidth=\textwidth]
    \begin{column}{0.49\textwidth}
      \centering
      \includegraphics[width=\linewidth,height=0.58\textheight,keepaspectratio]{#1}
      \vspace{0.04cm}
      
      {\scriptsize #2}
    \end{column}
    \begin{column}{0.49\textwidth}
      \centering
      \includegraphics[width=\linewidth,height=0.58\textheight,keepaspectratio]{#3}
      \vspace{0.04cm}
      
      {\scriptsize #4}
    \end{column}
  \end{columns}
}


% Thông tin dùng chung cho báo cáo và slide.
% Kiểm tra mã môn, lớp, học kỳ, năm học và giảng viên trước khi nộp.
\newcommand{\reporttitle}{BÁO CÁO ĐỒ ÁN: CLIPSeg}
\newcommand{\projecttopic}{Phân đoạn ảnh với text và visual prompt}
\newcommand{\groupname}{Cá nhân}
\newcommand{\coursename}{Ứng dụng xử lý ảnh và video số}
\newcommand{\courseschedule}{[Điền lịch học]}
\newcommand{\coursecode}{[Điền mã môn]}
\newcommand{\classname}{[Điền lớp]}
\newcommand{\semester}{[Điền học kỳ]}
\newcommand{\academicyear}{[Điền năm học]}
\newcommand{\studentone}{Lưu Huy Minh Quang}
\newcommand{\studentoneid}{23127016}
\newcommand{\studentoneemail}{lhmquang232@clc.fitus.edu.vn}
\newcommand{\studentonemajor}{Thị giác máy tính}
\newcommand{\advisorone}{[Điền giảng viên phụ trách]}
\newcommand{\reportplace}{Thành phố Hồ Chí Minh}
\newcommand{\reportdate}{tháng 10 năm 2026}

\title[CLIPSeg]{CLIPSeg: phân đoạn ảnh theo prompt}
\subtitle{Khảo sát mô hình và thực nghiệm trên tập nhỏ}
\author[\studentone]{\studentone{} -- \studentoneid}
\institute[HCMUS]{Trường Đại học Khoa học Tự nhiên, ĐHQG-HCM}
\date{}

\newcommand{\figureplaceholder}[2][0.45\textheight]{%
  \fbox{\parbox[c][#1][c]{0.86\linewidth}{\centering\small #2}}%
}

\begin{document}
\begin{frame}[plain]
\centering
\vspace{0.7cm}
\makebox[\textwidth][c]{%
\begin{beamercolorbox}[wd=0.88\textwidth,rounded=true,sep=0.45cm,center]{frametitle}
{\LARGE\bfseries CLIPSeg: phân đoạn ảnh theo prompt\par}
\vspace{0.22cm}
{\large Khảo sát mô hình và thực nghiệm trên tập nhỏ\par}
\end{beamercolorbox}}
\vspace{0.45cm}
{\large \studentone{} -- \studentoneid\par}
\vspace{0.2cm}
{\normalsize \coursename\par}
\vspace{0.15cm}
{\normalsize \reportdate\par}
\end{frame}

\begin{frame}{Nội dung trình bày}
\begin{enumerate}
\item Bài toán và phạm vi đồ án
\item Kiến trúc CLIPSeg và các loại prompt
\item Dữ liệu và protocol đánh giá
\item Thực nghiệm E1--E3 và phân tích lỗi
\item Demo, kết luận và hạn chế
\end{enumerate}
\end{frame}

\section{Giới thiệu}
\begin{frame}{Bài toán phân đoạn theo prompt}
\begin{columns}[T,totalwidth=\textwidth]
\begin{column}{0.48\textwidth}
\begin{itemize}
\item Đầu vào: ảnh query và text hoặc ảnh support.
\item Đầu ra: mask nhị phân cho mục tiêu.
\item GT dùng để đánh giá chất lượng dự đoán.
\end{itemize}
\end{column}
\begin{column}{0.48\textwidth}
\centering
\figureplaceholder[0.40\textheight]{Chèn ảnh query, prompt và mask mục tiêu}
\end{column}
\end{columns}
\end{frame}

\begin{frame}{Mục tiêu và phạm vi}
\begin{itemize}
\item Mô hình chính: \texttt{CIDAS/clipseg-rd64-refined}.
\item Triển khai inference với checkpoint có sẵn.
\item E1 đo baseline, E2 so text/visual, E3 khảo sát template.
\item Kết luận dựa trên số đo và ví dụ lỗi thực tế.
\end{itemize}
\end{frame}

\begin{frame}{Các phương pháp liên quan}
\begin{table}
\centering
\begin{tabularx}{\textwidth}{@{}l >{\raggedright\arraybackslash}X >{\raggedright\arraybackslash}X@{}}
\toprule
Phương pháp & Prompt & Đầu ra / điểm so sánh \\
\midrule
CLIPSeg & Text, ảnh mẫu & Mask nhị phân \\
{[Bổ sung 2]} & [Bổ sung] & [Bổ sung] \\
{[Bổ sung 3]} & [Bổ sung] & [Bổ sung] \\
\bottomrule
\end{tabularx}
\end{table}
\smallnote{Bổ sung nguồn đã kiểm chứng cho từng phương pháp.}
\end{frame}

\section{Phương pháp}
\begin{frame}{Kiến trúc CLIPSeg}
\centering
\figureplaceholder[0.52\textheight]{Chèn sơ đồ CLIP, vector điều kiện, FiLM và decoder}
\smallnote{Nguồn: L\"uddecke và Ecker, CVPR 2022.}
\end{frame}

\begin{frame}{Conditioning và decoder}
\begin{columns}[T,totalwidth=\textwidth]
\begin{column}{0.49\textwidth}
\begin{block}{Đặc trưng ảnh}
[Giải thích đặc trưng đa mức từ backbone CLIP của ảnh query.]
\end{block}
\end{column}
\begin{column}{0.49\textwidth}
\begin{block}{Điều kiện theo prompt}
[Giải thích vector text/visual, FiLM và cách decoder tạo logits.]
\end{block}
\end{column}
\end{columns}
\smallnote{Đối chiếu cấu hình checkpoint với biến thể paper khi nêu số tham số.}
\end{frame}

\begin{frame}{Text prompt và visual prompt}
\begin{columns}[T,totalwidth=\textwidth]
\begin{column}{0.49\textwidth}
\begin{block}{Text prompt}
[Chèn ví dụ câu nhắc và vùng tương ứng trên ảnh query.]
\end{block}
\end{column}
\begin{column}{0.49\textwidth}
\begin{block}{Visual prompt}
[Chèn ảnh support khác query và giải thích crop / giảm ảnh hưởng nền.]
\end{block}
\end{column}
\end{columns}
\end{frame}

\section{Thực nghiệm}
\begin{frame}{Dữ liệu và thiết kế đánh giá}
\begin{itemize}
\item Tập dữ liệu: [điền nguồn, split, số ảnh và query].
\item Tách calibration/evaluation theo image\_id, seed 42.
\item E2 dùng support từ ảnh khác query.
\item Các biến thể cùng ảnh, GT và cấu hình đánh giá.
\end{itemize}
\end{frame}

\begin{frame}{Hậu xử lý và metric}
\begin{itemize}
\item Nội suy logits về kích thước ảnh gốc, rồi sigmoid.
\item Ngưỡng mặc định 0,3, ngưỡng bổ sung chọn trên calibration.
\item Báo mean IoU theo query, cumulative IoU và AP pixel.
\end{itemize}
\[
\operatorname{meanIoU}=\frac{1}{N}\sum_i\frac{I_i}{U_i},\qquad
\operatorname{cIoU}=\frac{\sum_i I_i}{\sum_i U_i}.
\]
\end{frame}

\begin{frame}{E1: Baseline text}
\begin{table}
\centering
\begin{tabular}{@{}lcccc@{}}
\toprule
Cấu hình & $N$ & mean IoU & cIoU & AP \\
\midrule
Ngưỡng 0,3 & --- & --- & --- & --- \\
Ngưỡng calibration & --- & --- & --- & --- \\
\bottomrule
\end{tabular}
\end{table}
\smallnote{Điền số tự đo, split và ngưỡng đã khóa. Bổ sung một nhận xét từ kết quả.}
\end{frame}

\begin{frame}{E2: Text so với visual prompt}
\begin{table}
\centering
\begin{tabular}{@{}lccc@{}}
\toprule
Prompt & mean IoU & cIoU & AP \\
\midrule
Text & --- & --- & --- \\
Support crop & --- & --- & --- \\
Support giảm nền & --- & --- & --- \\
\bottomrule
\end{tabular}
\end{table}
\smallnote{Điền $N$, nguồn support và một trường hợp text/visual khác nhau.}
\end{frame}

\begin{frame}{E3: Ảnh hưởng của template}
\begin{table}
\centering
\begin{tabular}{@{}lcc@{}}
\toprule
Template & mean IoU & AP \\
\midrule
\texttt{\{phrase\}} & --- & --- \\
\texttt{a photo of \{phrase\}} & --- & --- \\
\texttt{the \{phrase\}} & --- & --- \\
\texttt{\{phrase\} in the image} & --- & --- \\
\bottomrule
\end{tabular}
\end{table}
\smallnote{Thêm kết luận dựa trên các query đã khóa.}
\end{frame}

\begin{frame}{Kiểm tra inference và demo}
\twopic{../manuscript/figures/demo_text.png}{Text: ``a spoon''}{../manuscript/figures/demo_visual.png}{Visual: crop chứa thìa}
\smallnote{Ví dụ kiểm tra API từ ảnh repo tác giả. Visual dùng crop cùng ảnh. Chưa có GT để đo chất lượng, không phải kết quả E2.}
\end{frame}

\section{Kết luận}
\begin{frame}{Kết luận và hạn chế}
\begin{itemize}
\item {[Điền phát hiện chính từ E1--E3.]}
\item {[Nêu giới hạn tập con, checkpoint và protocol.]}
\item {[Chèn đường dẫn demo/video và repo khi hoàn tất.]}
\end{itemize}
\smallnote{Paper chính: L\"uddecke và Ecker, \textit{Image Segmentation Using Text and Image Prompts}, CVPR 2022.}
\end{frame}
\end{document}
